% titel: Erdsatellit: Kreisbahn, Umlaufzeit, geostationaer
% skript: Kap. Gravitation, Satellitenbewegung

Ein Satellit der Masse $m\ll M_E$ kreist auf einer (näherungsweise)
Kreisbahn mit Radius $R$ um die Erde. Die Gravitationskraft liefert die
Zentripetalkraft; daraus folgen Bahngeschwindigkeit und Umlaufzeit allein aus
Bahnradius, Erdmasse und Gravitationskonstante.

\begin{werkzeugkasten}[Werkzeuge]
\textbf{Zentripetalkraft aus Gravitation:}
\begin{equation}
  \frac{m v^2}{R} = G\,\frac{M_E\,m}{R^2}
  \;\Rightarrow\;
  v = \sqrt{\frac{G M_E}{R}}\;.
\end{equation}

\textbf{Umlaufzeit} (Kreisbahn, $2\pi R$ pro Umlauf):
\begin{equation}
  T = \frac{2\pi R}{v} = 2\pi\sqrt{\frac{R^3}{G M_E}}\;.
\end{equation}

\textbf{Bahnhöhe über Erdoberfläche:} $R = R_E + h$.

\textbf{Geostationäre Bahn:} $T$ gleich der Umlaufdauer der Erde
($T=$\,\unit{24 h}); der Satellit steht dann über einem festen Punkt des
Äquators. Nur über dem Äquator ist die Bahnebene mit der
Drehachse verträglich.

\textbf{Mechanische Energie auf Kreisbahn:}
\begin{equation}
  E_\mathrm{mech} = E_\mathrm{kin}+E_\mathrm{pot}
  = \frac12 m v^2 - G\frac{M_E m}{R}
  = -G\frac{M_E m}{2R}\;\;(\text{mit } v^2=GM_E/R)\;.
\end{equation}

\textbf{Konstante:} $G=\qty{6.674e-11}{m^3/(kg\,s^2)}$,
$M_E=\qty{5.974e24}{kg}$, $R_E=\qty{6378}{km}$.
\end{werkzeugkasten}

\begin{werkzeugkasten}[Vorgehen in allgemeiner Form]
$a)$ Allgemeinen Ausdruck für $T$ herleiten. $b)$ Für die ISS Höhe
$h=\qty{400}{km}$: $R=R_E+h$ einsetzen. $c)$ Geostationär: $T$ fixieren und
nach $R$ auflösen. $d)$ Reibungsabnahme der Gesamtenergie $E_\mathrm{mech}$
mit dem Radius $R$ verknüpfen ($E_\mathrm{mech}\propto -1/R$) und
Schlussfolgerung über $R$ und $T$.
\end{werkzeugkasten}

\begin{achtung}[Typische Fallen]
\begin{itemize}
  \item $R$ ist der Abstand zum Erdmittelpunkt, nicht die Höhe über der
       Oberfläche.
  \item $E_\mathrm{mech}<0$: eine kleinere (negativere) mechanische Energie
       bedeutet kleineren Bahnradius.
\end{itemize}
\end{achtung}
